\documentclass[12pt, a4paper]{article}
\usepackage[utf8]{inputenc}
\usepackage[ngerman]{babel}
\usepackage{amsmath, amssymb}
\usepackage{geometry}
\geometry{margin=2.5cm}
\usepackage{parskip}

\title{\textbf{Aufgabe 7 -- Lösung} \\ Methoden I: Mathematik, 2026S}
\author{}
\date{}

\begin{document}
\maketitle

\section*{Angabe}
Multiplizieren Sie $(3x - 2)^6$ aus. Bestimmen Sie die Binomialkoeffizienten mit dem Pascalschen Dreieck.

\section*{Der Binomische Lehrsatz}

\[
    (a + b)^n = \sum_{k=0}^{n} \binom{n}{k} a^{n-k} \cdot b^k
\]

Hier: $a = 3x$, $b = -2$, $n = 6$.

\section*{Schritt 1: Pascalsches Dreieck (bis Zeile 6)}

\begin{center}
\begin{tabular}{rccccccccccccc}
    $n=0$: & & & & & & & 1 \\
    $n=1$: & & & & & & 1 & & 1 \\
    $n=2$: & & & & & 1 & & 2 & & 1 \\
    $n=3$: & & & & 1 & & 3 & & 3 & & 1 \\
    $n=4$: & & & 1 & & 4 & & 6 & & 4 & & 1 \\
    $n=5$: & & 1 & & 5 & & 10 & & 10 & & 5 & & 1 \\
    $n=6$: & 1 & & 6 & & 15 & & 20 & & 15 & & 6 & & 1
\end{tabular}
\end{center}

\textbf{Wie funktioniert das?} Jede Zahl ist die Summe der zwei Zahlen direkt darüber. Z.B. $15 = 5 + 10$.

Also: $\binom{6}{0} = 1$, $\binom{6}{1} = 6$, $\binom{6}{2} = 15$, $\binom{6}{3} = 20$, $\binom{6}{4} = 15$, $\binom{6}{5} = 6$, $\binom{6}{6} = 1$.

\section*{Schritt 2: Alle Terme berechnen}

\begin{align*}
    (3x - 2)^6 &= \sum_{k=0}^{6} \binom{6}{k} (3x)^{6-k} \cdot (-2)^k
\end{align*}

\begin{tabular}{c|c|c|c|c}
    $k$ & $\binom{6}{k}$ & $(3x)^{6-k}$ & $(-2)^k$ & Produkt \\
    \hline
    0 & 1 & $(3x)^6 = 729x^6$ & $1$ & $729x^6$ \\
    1 & 6 & $(3x)^5 = 243x^5$ & $-2$ & $6 \cdot 243 \cdot (-2) \cdot x^5 = -2916x^5$ \\
    2 & 15 & $(3x)^4 = 81x^4$ & $4$ & $15 \cdot 81 \cdot 4 \cdot x^4 = 4860x^4$ \\
    3 & 20 & $(3x)^3 = 27x^3$ & $-8$ & $20 \cdot 27 \cdot (-8) \cdot x^3 = -4320x^3$ \\
    4 & 15 & $(3x)^2 = 9x^2$ & $16$ & $15 \cdot 9 \cdot 16 \cdot x^2 = 2160x^2$ \\
    5 & 6 & $(3x)^1 = 3x$ & $-32$ & $6 \cdot 3 \cdot (-32) \cdot x = -576x$ \\
    6 & 1 & $(3x)^0 = 1$ & $64$ & $64$
\end{tabular}

\textbf{Hinweis:} $(-2)^k$ wechselt das Vorzeichen: positiv wenn $k$ gerade, negativ wenn $k$ ungerade.

\section*{Ergebnis}

\[
    \boxed{(3x-2)^6 = 729x^6 - 2916x^5 + 4860x^4 - 4320x^3 + 2160x^2 - 576x + 64}
\]

\end{document}
